\begin{table}[htbp]\centering
\caption{Direct expected-cost comparisons under a predeclared decision tolerance}\label{tab:direct47}
\begin{tabular}{lrrrr}\toprule
Contrast & Groups & Signed & Within $1/1024$ & Largest endpoint\\\midrule
Witness $-$ Cone nearest & 40 & 0 & 40 & 0.0001161 \\
Witness $-$ Spline & 40 & 0 & 40 & 0.0001975 \\
Cone nearest $-$ Spline & 40 & 0 & 40 & 0.0002078 \\
\bottomrule\end{tabular}
\par\smallskip\begin{minipage}{0.98\linewidth}\footnotesize Signed counts exclude intervals containing zero. ``Within'' requires the entire interval inside $[-1/1024,1/1024]$. Largest endpoint is the largest absolute endpoint over the 40 groups. All 120 intervals use one $0.01$ family error account, conditional on the frozen policies and independent-bin model. Each group uses 131,072 common paths.\end{minipage}\end{table}

Every direct interval contains zero, so the experiment identifies no signed policy ranking. Every interval is also contained in the predeclared decision band. It therefore certifies two-sided expected-cost proximity for the specified initial laws at that tolerance, rather than inferring equivalence from a failure to reject equality. The uniform initial-state law is a distributional conclusion, not an all-state assertion. The individual initial states are separate declared estimands.
